% GBC: global tensor boundary, original theta return, and relative cone.
\section{The all-prime mixed boundary in the actual theta comparison complex}
\label{sec:global-boundary-cone-lift}

\paragraph{Sources and the original arrows.}
The all-prime source and boundary are BW1--24 of
\nolinkurl{ROOT_BOUNDARY_WEIGHT_HOMOLOGY.tex}, in the shared
programme directory \texttt{rh\_joint\_strategy\_20260923}; the
complete proof was read for this calculation. Its original
public source is
\href{https://github.com/KokunoYumeto/zeta-function-research-reader/blob/3c78cfbd9e194d35117b33277c421c6564ce3dbf/workbenches/splitzero-tandem/continuations/20260923-full-prime-support/031/01_REBUILT_MATHEMATICS.tex}{\emph{Full prime-support reconstruction}, AB1--5, AB20--24}.
The original theta and cohomological trace source is Alain Connes,
Caterina Consani and Matilde Marcolli,
\href{https://arxiv.org/abs/math/0703392v1}{\emph{The Weil proof and
the geometry of the adeles class space}}, locators
\texttt{rangecVdef}, \texttt{varthetaH1V}, \texttt{vanishV} and
\texttt{tracepairing}. SWQ1--28 and GDT1--27 retain its exact
test algebra, closed quotient, full tensor action and duality.
The constructions below use the full rational group, including
its sign, and the original tensor before any endpoint quotient.

\subsection{The original Schwartz source and both endpoint amplitudes}

Write $G=\mathbb Q^\times$ and retain the original Schwartz space
and rational action
\begin{equation}\tag{GBC1}
 T=\mathcal S(\mathbb A_{\mathbb Q}),\quad
 (U_g\xi)(x)=\xi(g^{-1}x),\quad
 E\xi(x)=\sum_{q\in\mathbb Q^\times}\xi(qx).
\end{equation}
Let $\mathscr B$ be the entire original weighted Bruhat test
algebra on $C=\mathbb A_{\mathbb Q}^\times/G$, and let
$V=E(T_0)$, where $T_0=\{\xi:\xi(0)=0,\int\xi=0\}$.
Set $W=\overline V$ in that original topology and $H=\mathscr B/W$.
No even projection on individual tensor factors is made.

Use the following actual even real Schwartz functions only to
specify an endpoint section; the source $T$ itself is unrestricted:
\begin{equation}\tag{GBC2}
 \eta_0(x)=(1-2\pi x^2)e^{-\pi x^2},\qquad
 \eta_1(x)=2\pi x^2e^{-\pi x^2},\qquad
 \xi_i=\eta_i\prod_p\mathbf1_{\mathbb Z_p},\quad \theta_i=E\xi_i.
\end{equation}
The Gaussian integral and its derivative give
$(\xi_0(0),\int\xi_0)=(1,0)$ and
$(\xi_1(0),\int\xi_1)=(0,1)$, with every finite Haar factor
equal to one. The original theta sums are
$\theta_i(y)=2\sum_{n\ge1}\eta_i(ny)$.
Poisson summation, with its zero term retained, gives
$\theta_0(y)=-1+o(y^N)$ and
$\theta_1(y)=y^{-1}+o(y^N)$ as $y\downarrow0$, for every
$N$, while both decrease to every order at infinity.
For example the Fourier transform interchanges $\eta_0,\eta_1$;
the two exact equations are
$\theta_0(y)=y^{-1}\theta_1(y^{-1})-1$ and
$\theta_1(y)=y^{-1}\theta_0(y^{-1})+y^{-1}$.
Thus the entire endpoint-enlarged target and map are
\begin{equation}\tag{GBC3}
 A=\mathscr B\oplus\Theta,\quad
 \Theta=\mathbb C\theta_0\oplus\mathbb C\theta_1,\quad
 \pi(f+c_0\theta_0+c_1\theta_1)=(c_0,c_1),\quad
 q_A(f+c_0\theta_0+c_1\theta_1)=[f]\in H.
\end{equation}
Directness follows from the two different displayed tails at zero.
For arbitrary $\xi\in T$ put
$\xi'=\xi-\xi(0)\xi_0-(\int\xi)\xi_1$.
Its two moments vanish, so $E\xi'\in V\subset\mathscr B$.
Consequently, with no removed endpoint in the equality,
\begin{equation}\tag{GBC4}
 E\xi=E\xi'+\xi(0)\theta_0+(\int\xi)\theta_1,
 \quad \pi E\xi=(\xi(0),\int\xi),\quad q_AE\xi=0,
 \quad E(T)=V\oplus\Theta.
\end{equation}
The rational action on $A$ is trivial, since its functions live
on $C$. Reindexing the full rational sum proves $EU_g=E$.
This includes $g=-1$ and retains the full original sign action
on the source. The actual idele dilation remains equivariant:
it changes the moment pair by $(c_0,c_1)\mapsto(c_0,|a|c_1)$.
The difference between a dilated endpoint section and the
corresponding scaled endpoint section is an element of $V$,
by applying GBC4 to their moment-zero difference. Thus $q_A$
intertwines the original idele action on $H$.

\subsection{The all-prime tensor square of maps, before coinvariants}

Use $\mathcal B_{\partial}=\operatorname{LC}
(\{0,1\}^{\{\infty\}\sqcup\mathcal P},\mathbb C)$ and the
actual map $b:T\to\mathcal B_{\partial}$ of BW17, obtained
by tensoring evaluation and additive integration at every place.
Let $\epsilon:\mathcal B_{\partial}\to\mathbb C^2$ evaluate
at the two constant support patterns. Its formula is
$\epsilon b\xi=(\xi(0),\int\xi)$, including all additive
Haar factors. For any fixed $r\ge1$ retain all the algebraic tensors
\begin{equation}\tag{GBC5}
 \begin{gathered}
 T_r=T^{\otimes r},\quad \mathcal B_r=\mathcal B_{\partial}^{\otimes r},
 \quad A_r=A^{\otimes r},\quad P_r=(\mathbb C^2)^{\otimes r},\\
 b_r=b^{\otimes r},\ E_r=E^{\otimes r},\
 \pi_r=\pi^{\otimes r},\ \epsilon_r=\epsilon^{\otimes r},\
 q_r=q_A^{\otimes r},\\
 J_r=\ker b_r,\quad I_r=\ker\pi_r
     =\sum_i A^{\otimes(i-1)}\otimes\mathscr B
                          \otimes A^{\otimes(r-i)},\\
 \pi_r E_r=\epsilon_r b_r,\qquad
 E_r(J_r)\subset I_r,\qquad q_rE_r=0.
 \end{gathered}
\end{equation}
Both $b_r$ and $\pi_r$ are onto: BW17 gives an actual section
of $b$; GBC2--3 gives one of $\pi$; tensor their sections.
The tensor kernel formula is GFT9, also proved by expanding
the displayed direct sum in GBC3.

The group $G$ acts diagonally on the original $T_r$ and
$\mathcal B_r$; it acts trivially on $A_r,P_r,H^{\otimes r}$.
All four arrows in the square are equivariant. In particular
\begin{equation}\tag{GBC6}
 E_r(U_g^{\otimes r}-1)t=0\quad(t\in T_r,g\in G).
\end{equation}
This is a literal equality of entire theta functions. Thus the
degree-zero theta return of BW22 is zero before projection to
$H^{\otimes r}$. It does not identify the source boundary class
with zero in its actual rational homology.

\subsection{A relative comparison retaining the complete chain and its lift}

Use the full bar chain complex $C_n(G,M)=M\otimes\mathbb C[G^n]$
for every original $G$ module $M$, with finite linear combinations.
Write a generator as $(m;g_1,\ldots,g_n)$ and retain the differential
\begin{equation}\tag{GBC7}
 \begin{split}
 d(m;g_1,\ldots,g_n)={}&(g_1m;g_2,\ldots,g_n)\\
 &+\sum_{i=1}^{n-1}(-1)^i
       (m;g_1,\ldots,g_ig_{i+1},\ldots,g_n)\\
 &+(-1)^n(m;g_1,\ldots,g_{n-1}).
 \end{split}
\end{equation}
The degree-zero differential is zero. Here $G$ is abelian, so
the action term has exactly this order convention. In $d^2$,
merging two adjacent pairs in the two possible orders cancels
with opposite signs, including the triple merge by associativity;
removing a last entry cancels its opposite-order term; and the
two first action terms cancel because $U_gU_h=U_{gh}$ and
$G$ is abelian. This proves $d^2=0$. Applying an equivariant
linear map to the coefficient is a chain map, as each term shows.
This uses every rational sign and every prime. The prime Koszul
chains of BW15 enter through alternating bar symbols; in degree
one $m\theta_p$ is $(m;p)$ and its differential is $(U_p-1)m$.

For a chain map $F:C_\bullet\to D_\bullet$ use the exact convention
\begin{equation}\tag{GBC8}
 \operatorname{Cone}(F)_n=D_n\oplus C_{n-1},\qquad
 d(a,t)=(d_Da+Ft,-d_Ct).
\end{equation}
Its square is zero because $d_DF=Fd_C$. Form these actual complexes:
\begin{equation}\tag{GBC9}
 \begin{gathered}
 \mathfrak C_r=\operatorname{Cone}(C_\bullet(G,T_r)
                         \xrightarrow{E_r}C_\bullet(G,A_r)),\\
 \mathfrak P_r=\operatorname{Cone}(C_\bullet(G,\mathcal B_r)
                         \xrightarrow{\epsilon_r}C_\bullet(G,P_r)),\\
 \mathfrak R_r=\operatorname{Cone}(C_\bullet(G,J_r)
                         \xrightarrow{E_r}C_\bullet(G,I_r)).
 \end{gathered}
\end{equation}
The original square GBC5 gives a short exact sequence of complexes
\begin{equation}\tag{GBC10}
 0\longrightarrow\mathfrak R_r\longrightarrow\mathfrak C_r
      \xrightarrow{\Gamma_r}\mathfrak P_r\longrightarrow0,
 \qquad \Gamma_r(a,t)=(\pi_ra,b_rt).
\end{equation}
It is a chain map because $\pi_rE_r=\epsilon_rb_r$.
Surjectivity in each degree follows from the two actual surjections
in GBC5, applied to each finite bar coefficient. Its kernel is
exactly $C_n(G,I_r)\oplus C_{n-1}(G,J_r)$ with the differential
GBC8. This proves every assertion of exactness without discarding
any mixed support pattern or tensor sign.

In particular take a genuine cycle $c\in C_j(G,\mathcal B_r)$
with $\epsilon_rc=0$ as a chain, and choose an original lift
$t\in C_j(G,T_r)$ with $b_rt=c$. Its cone cycle is $(0,c)$
in degree $j+1$ of $\mathfrak P_r$. The exact connecting map is
\begin{equation}\tag{GBC11}
 \partial_{\rm cone}[(0,c)]=[(E_rt,-dt)]\in H_j(\mathfrak R_r).
\end{equation}
Indeed $\pi_rE_rt=\epsilon_rc=0$ and $b_rdt=dc=0$, so
the pair belongs to the displayed complex. Its differential is
$(dE_rt-E_rdt,d^2t)=0$. A new lift $t+u$, with $u$ in
$C_j(G,J_r)$, changes it by $(E_ru,-du)=d(0,u)$ in
$\mathfrak R_r$. For a general change by a cone boundary, lift
that bounding cone chain to $\mathfrak C_r$ using the degreewise
surjection; its differential lifts the change and its second
differential is zero. Any difference from that chosen lift is
in $\mathfrak R_r$ and contributes a boundary there. This proves
the full homology well-definedness of the connecting map.

For the actual degree-one BW22 input $c=(b_0;p)$, with
$(D_p^{(r)}-1)b_0=0$ and $\epsilon_rb_0=0$, it becomes
\begin{equation}\tag{GBC12}
 \left[(\,(E_rt_0;p),\ (1-U_p^{\otimes r})t_0\,)\right],
               \qquad b_rt_0=b_0.
\end{equation}
The second component is the negative of BW22 with the cone sign
retained. Its theta sum alone is zero by GBC6. The first component
retains the original lifted theta amplitude and is essential to the
relative chain. No assertion that this cone class is nonzero has
been inferred merely from a nonzero boundary coefficient.

\subsection{Its exact relation to the interior tensor and the remaining complex}

The canonical original interior receiver is the chain map
\begin{equation}\tag{GBC13}
 F_r:\mathfrak R_r\longrightarrow C_\bullet(G,H^{\otimes r}),
                    \qquad (a,j)\longmapsto q_ra.
\end{equation}
It commutes with the differential because $q_rE_r=0$ and $q_r$
is equivariant. It is onto in each degree: $I_r$ contains
$\mathscr B^{\otimes r}$, whose quotient onto $H^{\otimes r}$
is onto. For the entire mixed connecting image there is the exact
chain-level equality
\begin{equation}\tag{GBC14}
 F_r(E_rt,-dt)=q_rE_rt=0.
\end{equation}
Thus this particular full theta comparison does not send the
mixed BW connecting classes to the nonzero balanced primary
tensors of GDT26--27. This is the proved value of the actual map,
not a conclusion that the spaces admit no other relation.

Its full kernel is a constructed comparison complex. Put
\begin{equation}\tag{GBC15}
 \begin{gathered}
 K_r=I_r\cap\ker q_r,\qquad
 \mathfrak K_r=\operatorname{Cone}(C_\bullet(G,J_r)
                          \xrightarrow{E_r}C_\bullet(G,K_r)),\\
 0\longrightarrow\mathfrak K_r\longrightarrow\mathfrak R_r
       \xrightarrow{F_r}C_\bullet(G,H^{\otimes r})\longrightarrow0.
 \end{gathered}
\end{equation}
The inclusion $E_r(J_r)\subset K_r$ follows from both equations
in GBC5. Computing the kernel of GBC13 in each degree gives
exactly the first two lines of GBC15, including their differential.
The all-prime connecting representative GBC11 is already a cycle
in this kernel complex, not just a class whose image happens to
vanish after passing to homology.

There is also an explicit entire description of its coefficient
space, retaining every tensor sector. Expand GBC3 with the
original order of factors; for $S\subset\{1,\ldots,r\}$ let
$A_S=\bigotimes_i M_i$, where $M_i=\mathscr B$ for $i\in S$
and $M_i=\Theta$ otherwise. Then
\begin{equation}\tag{GBC16}
 \begin{gathered}
 A_r=\bigoplus_{S\subset[r]}A_S,\qquad
 I_r=\bigoplus_{\varnothing\ne S\subset[r]}A_S,\\
 K_r=\left(\bigoplus_{\varnothing\ne S\subsetneq[r]}A_S\right)
       \oplus\sum_{i=1}^r
         \mathscr B^{\otimes(i-1)}\otimes W
                         \otimes\mathscr B^{\otimes(r-i)}.
 \end{gathered}
\end{equation}
The first equality is the finite distributive expansion of a
vector-space direct sum, preserving each original position.
The endpoint map is an isomorphism on the all-$\Theta$ sector
and zero on every other sector. The map $q_r$ is zero on every
sector containing $\Theta$; on the all-$\mathscr B$ sector its
kernel is precisely the tensor quotient kernel GFT9. These facts
prove all lines of GBC16. In particular for $r=2$ the mixed
spaces $\mathscr B\otimes\Theta$ and $\Theta\otimes\mathscr B$
both remain, along with the full closed-image tensor kernel.
The splitting into $\Theta$ sectors is an explicit vector-space
section, not a claim that each sector is invariant under idele
dilation. The subspaces $I_r,K_r$ and all maps are invariant;
the section's extension terms are the original elements of $V$
computed after GBC4.

\subsection{The joint source kernel containing the entire mixed connecting image}

The preceding receiver can be made more exact by calculating its
original theta-image intersection. First,
\begin{equation}\tag{GBC17}
 E(\ker b)=V.
\end{equation}
The inclusion into $V$ holds because $b\xi=0$ implies both
global moments vanish. Conversely let $E\xi\in V$ with
$\xi\in T_0$. Then $b\xi\in\ker\epsilon$. The proved full
boundary relation BW6 at $r=1$ gives a finite expression
$b\xi=\sum_p(D_p-1)v_p$. Lift each coefficient $v_p$ by the
onto original map $b$ to $t_p\in T$. The explicit vector
$j=\xi-\sum_p(U_p-1)t_p$ belongs to $\ker b$, while
$Ej=E\xi$ by the full rational invariance GBC6. This proves
GBC17 with an actual antecedent. No individual sign projection
or infinite sum of group relations is used.

Put $U=E(T)=V\oplus\Theta$. Tensoring the exact source kernel
formula BW18 and then applying $E_r$ gives
\begin{equation}\tag{GBC18}
 \begin{split}
 E_r(J_r)
   &=\sum_i U^{\otimes(i-1)}\otimes V\otimes U^{\otimes(r-i)}\\
   &=U^{\otimes r}\cap I_r
    =E_r\bigl(\ker(\epsilon_rb_r)\bigr).
 \end{split}
\end{equation}
For the first equality use GBC17 in the factor belonging to
$\ker b$ and surjectivity of $E:T\to U$ in every other factor;
this proves both inclusions on arbitrary finite sums of tensors.
For the second, expand $U=V\oplus\Theta$ inside GBC16.
The endpoint map is injective on the all-$\Theta$ summand and
zero on every summand with a $V$ factor. Thus its kernel is
exactly the displayed sum. Finally a tensor $t$ has
$\epsilon_rb_rt=0$ exactly when $E_rt\in I_r$, by GBC5;
the first two equalities give the last one.

The resulting original source spaces and their entire exact
sequence are
\begin{equation}\tag{GBC19}
 \begin{gathered}
 K_r^E=\ker E_r\subset T_r,\qquad
 K_r^{b,E}=\ker b_r\cap\ker E_r,\qquad
 D_r^\epsilon=\ker\epsilon_r\subset\mathcal B_r,\\
 0\longrightarrow K_r^{b,E}\longrightarrow K_r^E
                  \xrightarrow{b_r}D_r^\epsilon\longrightarrow0.
 \end{gathered}
\end{equation}
To prove onto, take $d\in D_r^\epsilon$ and any original
lift $t$ with $b_rt=d$. Then $E_rt\in U^{\otimes r}\cap I_r$,
so GBC18 supplies $u\in J_r$ with $E_ru=E_rt$.
The vector $t'=t-u$ lies in $K_r^E$ and still satisfies
$b_rt'=d$. The kernel of this restricted map is exactly
$K_r^{b,E}$. This proves the full sequence, with all spaces
subspaces of the original Schwartz tensor. Equivariance follows
from the original equivariance of $b_r,E_r$.

Apply this coefficient construction to the finitely many bar
coefficients of a cycle $c\in C_j(G,D_r^\epsilon)$.
It gives $t'$ with $b_rt'=c$ and $E_rt'=0$.
Then $dt'\in C_{j-1}(G,K_r^{b,E})$, and the exact connecting
map of GBC19 and its relation to the relative cone are
\begin{equation}\tag{GBC20}
 \begin{gathered}
 H_j(G,D_r^\epsilon)\longrightarrow
              H_{j-1}(G,K_r^{b,E}),\qquad[c]\longmapsto[dt'],\\
 [(E_rt,-dt)]=[(0,-dt')]
                    \quad\hbox{in }H_j(\mathfrak R_r),\qquad t'=t-u.
 \end{gathered}
\end{equation}
The first map is independent of the lift since a change lies
in $C_j(G,K_r^{b,E})$ and changes the answer by its boundary.
For a change of $c$ by a boundary, lift the bounding chain and
use $d^2=0$, as in the proof of GBC11. For the second equality,
subtract the literal relative-cone boundary
$d(0,u)=(E_ru,-du)$ from $(E_rt,-dt)$.
The result is exactly $(0,-d(t-u))$ with the stated sign.
Thus the entire secondary class comes from the shifted chain
complex of the explicit joint source kernel; it has not been
declared zero. The full source extension GBC19 records which
boundary homology classes it reaches. This calculation proves
the exact relation rather than proposing an unidentified map
between spaces having the same modulus weight.

The group chains, cone differentials, exact sequences and tensor
receivers lift over the full $G_L(\mathbb C)$ by GFT3--15a:
retain one original label on each whole amplitude chain and apply
every displayed linear map without changing it. Tensor labels
use their complete meet; finite chain sums use their join.
A squared differential has amplitude zero and returns the same
label, including $e$ when it was top, rather than $\tau$.
There is no infinite lattice join and no independent-coordinate
label quotient hidden in these complexes. GBC10 and GBC15 have
the exact same-label kernel congruences proved in GFT10.

The idele action commutes with all original differentials and
arrows. Therefore the connecting maps respect its actual
characters on homology: the boundary character $|a|^m$ computed
in BW14 is retained on the reached relative class, including
its zero value when the class vanishes. In the modulus convention
$|a|^{w/2}$ this is weight $2m$. GBC14 proves that the displayed
interior receiver annihilates the whole mixed connecting image.
GBC15--16 constructs the precise relative object containing it,
with its full source chain, endpoint sectors and support data.
This calculation does not replace it by the two endpoint
coinvariants and does not assume an interior purity theorem.
